\section{无参考评估与 LLM 作为评估者}

\subsection{用 LLM 替代人工评估}

一个非常自然的想法是：既然人工评估太慢太贵，能不能用 GPT-4 等强大的 LLM 来代替人类做评估？

\begin{figure}[H]
\centering
\includegraphics[width=0.95\textwidth]{slides-images/slide-034.jpg}
\caption{LLM 作为评估者的工作流程：给定一条指令和两个模型的回答，让 GPT-4 判断哪个更好\protect\footnotemark}
\end{figure}
\footnotetext{Slides 第35页。}

\begin{importantbox}{LLM 评估的惊人发现}
使用 GPT-4 进行评估相比人工评估：
\begin{itemize}
    \item 快 \textbf{100 倍}
    \item 便宜 \textbf{100 倍}
    \item 与人类评估的一致性\textbf{高于}人类标注者之间的一致性！
\end{itemize}
最后一点看似矛盾，但原因在于：LLM 的\textbf{方差极低}（几乎总是给出一致的判断），而人类标注者的方差很高。即使 LLM 有一定的偏差（bias），其低方差使得总体一致性反而更高。
\end{importantbox}

\subsection{偏差-方差权衡}

\begin{figure}[H]
\centering
\includegraphics[width=0.95\textwidth]{slides-images/slide-036.jpg}
\caption{偏差-方差分析：人类标注者方差约 31--33\%、偏差为 0；GPT-4 方差极低但偏差约 32\%。总体而言，GPT-4 的低方差使其与人类的一致性反而更高\protect\footnotemark}
\end{figure}
\footnotetext{Slides 第37页。}

\begin{knowledgebox}{为什么 LLM 能比人类更``一致''}
人类的偏差（bias）按定义为 0（人类评估就是金标准），但方差（variance）很高——同一个人在不同时间、不同标注者对同一样本会给出不同评价。LLM 的偏差不为零（它有系统性的偏好），但方差极低——同样的输入几乎总是给出同样的输出。在偏差-方差权衡下，LLM 的低方差使其在整体一致性上表现出色。这也使得 LLM 评估对研究更加\textbf{友好}——结果可复现、不依赖特定标注者。
\end{knowledgebox}

\subsection{LLM 评估的偏差问题}

尽管 LLM 评估效果惊人，但需要警惕以下偏差：

\begin{figure}[H]
\centering
\includegraphics[width=0.95\textwidth]{slides-images/slide-037.jpg}
\caption{LLM 评估中的偏差：长度偏差（人类和模型都约 70\% 偏好更长的回答）、列表偏差（偏好包含列表的回答）和位置偏差（偏好特定位置的回答）\protect\footnotemark}
\end{figure}
\footnotetext{Slides 第38页。}

\begin{warningbox}{LLM 评估的三大偏差}
\begin{enumerate}
    \item \textbf{长度偏差}：人类和 LLM 都偏好更长的回答（约 70\% 的偏好率），即使长度与质量无关
    \item \textbf{格式偏差}：偏好包含列表、结构化格式的回答
    \item \textbf{位置偏差}：在并排比较中，回答的左右位置会影响判断（可通过随机化缓解）
    \item \textbf{自我偏好}（Self-bias）：GPT-4 评估时倾向于偏好 GPT-4 生成的回答——但这种偏好没有想象中那么严重
\end{enumerate}
\end{warningbox}

\subsection{AlpacaEval：高效的自动评估基准}

\begin{figure}[H]
\centering
\includegraphics[width=0.95\textwidth]{slides-images/slide-040.jpg}
\caption{AlpacaEval 工作原理：给定指令，比较待评估模型与基线模型的回答，由 GPT-4 给出偏好概率，经过长度去偏重加权后计算胜率\protect\footnotemark}
\end{figure}
\footnotetext{Slides 第41页。}

AlpacaEval 是 Yann Dubois 等人在开发 Alpaca 项目时为超参调优而开发的轻量级基准，后来发展成为广泛使用的评估工具：

\begin{itemize}
    \item 与 Chatbot Arena 排名的\textbf{相关性高达 98\%}
    \item 评估一个模型仅需约 \textbf{3 分钟}和 \textbf{10 美元}
    \item 工作原理：给定指令 → 待评估模型和基线模型各生成一个回答 → GPT-4 给出偏好概率 → 经过重加权后计算胜率
\end{itemize}

\begin{figure}[H]
\centering
\includegraphics[width=0.95\textwidth]{slides-images/slide-041.jpg}
\caption{系统级相关性：AlpacaEval（x 轴）与 Chatbot Arena（y 轴，金标准）高度相关；下方图表显示 MT-Bench、AlpacaEval 和 MMLU 与 Chatbot Arena 的相关系数\protect\footnotemark}
\end{figure}
\footnotetext{Slides 第42页。}

\subsubsection{长度去偏的必要性}

\begin{figure}[H]
\centering
\includegraphics[width=0.95\textwidth]{slides-images/slide-043.jpg}
\caption{长度偏差的影响：GPT-4 被提示``更详细回答''时胜率从 50\% 飙升至 64.3\%，被提示``更简洁''时降至 22.9\%——这不符合我们对基准测试的预期。经过长度去偏重加权后，胜率恢复稳定\protect\footnotemark}
\end{figure}
\footnotetext{Slides 第44页。}

\begin{importantbox}{为什么需要长度去偏}
如果仅通过调整 prompt 让模型回答更详细，胜率就能从 50\% 跳到 64.3\%——这完全不是真实能力的提升。AlpacaEval 通过长度去偏重加权解决了这个问题：如果回答更长，给予略低的偏好权重，使最终分数不受长度操控的影响。
\end{importantbox}

\subsubsection{自我偏好问题}

\begin{figure}[H]
\centering
\includegraphics[width=0.95\textwidth]{slides-images/slide-045.jpg}
\caption{自我偏好分析：不同评估模型（列）对不同被评模型（行）的排名。尽管每个模型确实给自己更高的分数（对角线），但整体排名在不同评估者间保持一致\protect\footnotemark}
\end{figure}
\footnotetext{Slides 第46页。}

自我偏好确实存在，但\textbf{没有想象中那么严重}。例如，Mistral 用自己评估时确实给自己更高的分，但仍然把 Claude 和 GPT-4 排在自己前面。不同评估模型产生的排名总体一致。

\subsection{本章小结}

LLM 作为评估者是近年来最重要的评估方法创新之一。GPT-4 等模型的评估比人工快 100 倍、便宜 100 倍，且与人类的一致性甚至高于人类之间的一致性（归因于低方差）。AlpacaEval 证明了 LLM 评估可以以极低的成本达到与 Chatbot Arena 98\% 的相关性。但需要注意长度偏差、格式偏差和自我偏好等问题，并通过去偏技术加以缓解。

%% ========== Part 6: LLM 评估的当前实践 ==========

